% ===== PRÉAMBULE CANONIQUE — à coller VERBATIM en tête de CHAQUE .tex =====
\documentclass[11pt,a4paper]{article}
\usepackage[utf8]{inputenc}\usepackage[T1]{fontenc}\usepackage{lmodern}
\usepackage[french]{babel}
\usepackage{amsmath,amssymb,mathtools}\usepackage{icomma}
\usepackage[version=4]{mhchem}          % \ce{...} : équations & formules
\usepackage{siunitx}\sisetup{locale=FR} % \qty{}{}, \si{}, \ang{}
\usepackage{chemfig}                    % \chemfig{...} : structures développées
\usepackage{xcolor}\usepackage{colortbl}
\usepackage{tikz}\usepackage{pgfplots}\pgfplotsset{compat=1.18}
\usepackage[most]{tcolorbox}\usepackage{enumitem}\usepackage{array,booktabs}
\usepackage[a4paper,margin=2.2cm]{geometry}
\usetikzlibrary{arrows.meta,calc,angles,quotes,positioning,patterns,backgrounds,shapes.geometric}
\definecolor{primary}{RGB}{222,49,99}\definecolor{primarydark}{RGB}{150,22,55}
\definecolor{defcol}{RGB}{0,114,178}\definecolor{methcol}{RGB}{52,92,148}
\definecolor{excol}{RGB}{0,135,90}\definecolor{attncol}{RGB}{200,40,55}
\tcbset{boxbase/.style={enhanced,breakable,boxrule=0.9pt,arc=2.5pt,
  left=3mm,right=3mm,top=2.3mm,bottom=2.3mm,fonttitle=\bfseries,coltitle=white}}
% --- boîtes TUTORIEL (noms EXACTS, ne pas renommer) ---
\newtcolorbox{cle}[1][]{boxbase,colback=defcol!6,colframe=defcol,
  title={\textsc{À retenir}\ifx\relax#1\relax\relax\else\textmd{ : \itshape #1}\fi}}
\newtcolorbox{methode}[1][]{boxbase,colback=methcol!7,colframe=methcol,sharp corners,
  title={\textsc{Méthode}\ifx\relax#1\relax\relax\else\textmd{ --- \itshape #1}\fi}}
\newtcolorbox{exemplebox}[1][]{boxbase,colback=excol!5,colframe=excol,
  title={\textsc{Exemple}\ifx\relax#1\relax\relax\else\textmd{ : \itshape #1}\fi}}
\newtcolorbox{piege}[1][]{boxbase,colback=attncol!6,colframe=attncol,
  title={\textsc{Piège}\ifx\relax#1\relax\relax\else\textmd{ : \itshape #1}\fi}}
\newtcolorbox{remarque}[1][]{boxbase,colback=black!4,colframe=black!45,
  title={\textsc{Remarque}\ifx\relax#1\relax\relax\else\textmd{ : \itshape #1}\fi}}
% --- boîtes EXERCICE / CORRIGÉ (noms EXACTS, auto-comptées) ---
\newtcolorbox[auto counter]{exo}[1][]{boxbase,colback=primary!4,colframe=primary,
  title={\textsc{Exercice}~\thetcbcounter\ifx\relax#1\relax\relax\else\textmd{ --- \itshape #1}\fi}}
\newtcolorbox[auto counter]{corrige}[1][]{boxbase,colback=excol!4,colframe=excol,
  title={\textsc{Corrigé}~\thetcbcounter\ifx\relax#1\relax\relax\else\textmd{ --- \itshape #1}\fi}}
\newcommand{\R}{\mathbb{R}}\newcommand{\N}{\mathbb{N}}\newcommand{\Z}{\mathbb{Z}}\newcommand{\ds}{\displaystyle}
% (un \usepackage{fancyhdr}+en-tête est possible ; il est ignoré côté web)
% =========================================================================
\begin{document}
\begin{exo}[reconnaître des isotopes]
Deux atomes sont des isotopes du même élément s'ils ont le \emph{même} $Z$ et des $A$
\emph{différents}. Pour chaque paire, dis si ce sont des isotopes, en justifiant.
\begin{enumerate}[label=\alph*)]
  \item $\ce{^{16}_{8}O}$ et $\ce{^{18}_{8}O}$
  \item $\ce{^{40}_{18}Ar}$ et $\ce{^{40}_{20}Ca}$
  \item $\ce{^{12}_{6}C}$ et $\ce{^{14}_{6}C}$
\end{enumerate}
\end{exo}

\begin{exo}[neutrons et chimie des isotopes]
On considère trois isotopes du magnésium : $\ce{^{24}_{12}Mg}$, $\ce{^{25}_{12}Mg}$ et
$\ce{^{26}_{12}Mg}$.
\begin{enumerate}[label=\alph*)]
  \item Donne le nombre de neutrons de chacun.
  \item Combien d'électrons possède chacun à l'état neutre ?
  \item Ces trois isotopes ont-ils les mêmes propriétés chimiques ? Justifie en une phrase.
\end{enumerate}
\end{exo}

\begin{exo}[masse atomique moyenne à deux isotopes]
Calcule la masse atomique relative moyenne $\overline{A}$ de chaque élément à partir de ses deux
isotopes et de leurs abondances.
\begin{enumerate}[label=\alph*)]
  \item chlore : $\ce{^{35}Cl}$ ($75\,\%$) et $\ce{^{37}Cl}$ ($25\,\%$) ;
  \item cuivre : $\ce{^{63}Cu}$ ($70\,\%$) et $\ce{^{65}Cu}$ ($30\,\%$) ;
  \item argent : $\ce{^{107}Ag}$ ($52\,\%$) et $\ce{^{109}Ag}$ ($48\,\%$).
\end{enumerate}
\end{exo}

\begin{exo}[retrouver une abondance manquante]
Les abondances de tous les isotopes d'un élément totalisent $100\,\%$. Détermine l'abondance manquante.
\begin{enumerate}[label=\alph*)]
  \item deux isotopes ; l'un a une abondance de $60\,\%$ ;
  \item trois isotopes d'abondances $78\,\%$, $10\,\%$ et $?$ ;
  \item quatre isotopes d'abondances $5{,}8\,\%$, $91{,}7\,\%$, $2{,}2\,\%$ et $?$.
\end{enumerate}
\end{exo}

\begin{exo}[l'unité de masse atomique]
On donne $\SI{1}{uma} \approx \SI{1,66054e-27}{\kilo\gram}$ et $N_{\!A} \approx \SI{6,022e23}{\per\mole}$.
\begin{enumerate}[label=\alph*)]
  \item Quelle est, en kilogrammes, la valeur d'une uma ?
  \item Quelle est, en uma, la masse approchée d'un atome de $\ce{^{23}Na}$ (un nucléon $\approx 1$ uma) ?
  \item Exprime $\SI{1}{uma}$ en grammes à l'aide de $N_{\!A}$.
\end{enumerate}
\end{exo}

\end{document}
